\documentclass[11pt,a4paper]{article}
\usepackage[utf8]{inputenc}
\usepackage{amsmath,amssymb,amsfonts,amsthm}
\usepackage{graphicx}
\usepackage{hyperref}

\title{Holographic Verification of Boundary-Forced Tiling Uniqueness and Absence of Phasons for the Spectre Monotile in Lean 4}
\author{James F. Walters}
\date{\today}

\begin{document}
\sloppy
\maketitle

\begin{abstract}
This paper documents an axiom-pure formal verification framework in the Lean 4 proof assistant establishing boundary-forced tiling uniqueness and the absence of phasons for the Spectre chiral aperiodic monotile ($\mathrm{Tile}(1,1)$). Rather than attempting an unconstrained global search over an exponential packing space, we employ a verified 3-Anchor-Lock priority state machine (\texttt{L2 > L1 > L4}) executed via computational reflection. We prove four core results with zero \texttt{sorry}s and zero custom axioms: (1) a deterministic 3-anchor-lock priority cascade completely and uniquely reduces the 71-tile Generation-2 $\Delta_2$ metatile boundary as well as all nine Tier-1 metatile boundaries ($\Gamma, \Delta, \Theta, \Lambda, \Xi, \Pi, \Sigma, \Phi, \Psi$); (2) every simple closed path constructed from edges of the unique interior tiling of $\Delta_2$ encloses a unique interior sub-patch; (3) the 71-step acyclic elimination order is mirrored directly in an interactive Progressive Web App (PWA); and (4) combining Combinatorial Gauss-Bonnet with an exhaustive 4-batch transfer-matrix curvature bound rules out minimal phasons and proves that full Spectre tilings admit no phason flips.
\end{abstract}

\section{Introduction}
\label{sec:intro}
Establishing boundary-forced tiling uniqueness and ruling out local phason flips are foundational problems in the theory of aperiodic order. Standard unconstrained backtracking search algorithms suffer from exponential state-space explosion and geometric frustration. For the Spectre chiral aperiodic monotile~\cite{smith2023chiral}, discovered by Smith, Myers, Kaplan, and Goodman-Strauss in 2023, the tile admits only chiral placements without mirror reflections. By fusing an exact cyclotomic lattice embedding with a deterministic 3-Anchor-Lock boundary reduction machine inside the Lean 4 proof assistant~\cite{moura2021lean4}, we prove that boundary turn constraints uniquely determine the interior tiling and prevent any local phason rearrangements.

\section{Algebraic Foundations: Exact Cyclotomic Lattice Embedding}
\label{sec:algebra}
To guarantee absolute safety against floating-point drift, we embed the tile geometry into the 4D cyclotomic integer ring $\mathbb{Z}[\zeta_{12}]$, where $\zeta_{12} = e^{i\pi/6}$. Every lattice point $p \in \mathbb{Z}[\zeta_{12}]$ is represented as a 4-tuple $(a, b, c, d) \in \mathbb{Z}^4$ relative to the integral basis $\{1, \zeta_{12}, \zeta_{12}^2, \zeta_{12}^3\}$. A counterclockwise rotation by $30^\circ$ is represented exactly by the linear operator:
\begin{equation}
R(a, b, c, d) = (-d, a, b + d, c)
\end{equation}
We formally prove in Lean 4 (\texttt{rot12\_identity} in \texttt{Bedrock.lean}) that applying this operator 12 times yields the identity mapping, providing an exact discrete representation of Euclidean rotations without numerical rounding.

\section{Combinatorial Geometry and the Unified Spatial Predicate}
\label{sec:sensors}
Physical intersections and boundary touches are unified using the exact coordinate predicate \texttt{discreteSegmentsIntersectOrTouch} in \texttt{Paths.lean}.
Given $\zeta_{12} = \frac{\sqrt{3}}{2} + \frac{1}{2}i$, a 4D lattice point $a + b\zeta_{12} + c\zeta_{12}^2 + d\zeta_{12}^3$ projects onto $\mathbb{Z}[\sqrt{3}] \times \mathbb{Z}[\sqrt{3}]$. Clearing the common denominator of 2 maps each coordinate to exact integer coefficients of $(1, \sqrt{3})$:
\begin{align}
X &= (2a + c) + b\sqrt{3} \\
Y &= (b + 2d) + c\sqrt{3}
\end{align}
Because the 2D cross-product determinant is homogeneous of degree 2, clearing the factor of 2 scales the determinant by 4 and preserves the exact orientation sign over $\mathbb{Z}[\sqrt{3}]$.

The canonical counterclockwise (CCW) perimeter of the Spectre monotile is encoded in \texttt{Monotile.lean} by the 14 inter-edge turns:
\begin{multline}
\texttt{spectreTurns} = [-60^\circ, 90^\circ, 60^\circ, -90^\circ, 60^\circ, 90^\circ, 60^\circ, \\
-90^\circ, 60^\circ, 0^\circ, 60^\circ, 90^\circ, -60^\circ, 90^\circ]
\end{multline}

\section{The 3-Anchor-Lock Priority Machine (\texttt{L2 > L1 > L4})}
\label{sec:cascade}
We formulate boundary-forced tiling uniqueness as a deterministic state machine (\texttt{Certificate.lean}) over closed loops of directed lattice edges. Analogous to Gonthier's formal proof of the Four Color Theorem~\cite{gonthier2008four}, our proof uses computational reflection inside the Lean 4 kernel.

Remarkably, the entire 71-tile Generation-2 $\Delta_2$ patch and all nine Tier-1 metatiles ($\Gamma, \Delta, \Theta, \Lambda, \Xi, \Pi, \Sigma, \Phi, \Psi$) are completely reduced using \textbf{strictly three anchor locks} evaluated in priority order (\texttt{selectPriorityLock}):
\begin{enumerate}
\item \textbf{Priority 1 --- Lock \#2 (\texttt{300049}, High Lock $[60^\circ, 90^\circ, 60^\circ]$):} CCW step signature $[2, 3, 2]$, spanning 4 consecutive tile edges (indices $4, 5, 6, 7$). Used in 55 of the 71 steps of $\Delta_2$.
\item \textbf{Priority 2 --- Lock \#1 (\texttt{300033}, Low Lock $[60^\circ, 0^\circ, -60^\circ]$):} CCW step signature $[2, 0, -2]$, spanning 3 consecutive tile edges (indices $8, 9, 10$). Used in 10 of the 71 steps of $\Delta_2$.
\item \textbf{Priority 3 --- Lock \#4 (\texttt{400110}, Long Lock $[90^\circ, -60^\circ, 90^\circ, 60^\circ]$):} CCW step signature $[3, -2, 3, 2]$, spanning 5 consecutive tile edges (indices $13, 0, 1, 2, 3$). Used in 6 of the 71 steps of $\Delta_2$.
\end{enumerate}

At each step, the validator converts the active boundary loop into a cyclically padded edge-turn stream (\texttt{loopToPaddedEdgeTurns}) in a single linear pass, selects the highest-priority lock present on the boundary, and executes \texttt{matchingAlignedTiles} across all 14 possible monotile edge alignments ($\texttt{tile\_edge\_idx} \in \{0, \dots, 13\}$) to verify that \textbf{uniquely one} placed tile covers the locked boundary segments. The state machine then annihilates shared boundary edges and exposes the reversed interior tile edges via \texttt{computeNextBoundaryEdges}, returning \texttt{Option PeelingState} so that any invalid transition immediately halts with \texttt{none}.

\section{Uniqueness of Interiors for Simple Closed Paths}
\label{sec:simplepaths}
In \texttt{SimplePaths.lean}, we extend uniqueness from the outer $\Delta_2$ perimeter to all simple closed paths (closed loops of tile edges visiting each vertex at most once) inside $\Delta_2$:
\begin{enumerate}
\item \textbf{Simple Loop Invariant (\texttt{delta2\_all\_cascade\_loops\_simple}):} The initial $\Delta_2$ perimeter and all 70 intermediate boundary loops in the peeling cascade are verified to be simple closed paths (\texttt{hasNoDuplicateNodes}).
\item \textbf{Acyclic 3-Lock Elimination Order (\texttt{delta2\_lock\_elimination\_sequence\_valid}):} Ordering the 71 interior tiles as $t_1, \dots, t_{71}$ by the peeling cascade, at every step $k \in \{1, \dots, 71\}$ the anchor-lock edges $L_k \subset \partial t_k$ uniquely force $[t_k]$ across all 14 monotile alignments, and no later tile $t_j$ ($j > k$) contains any edge of $L_k$ or its reversal.
\item \textbf{Universal Sub-Patch Uniqueness (\texttt{simple\_closed\_path\_unique\_interior}, \texttt{subpatch\_unique\_of\_same\_boundary}):} For any non-empty sub-patch $P \subseteq \{t_1, \dots, t_{71}\}$ bounded by a simple closed path $C = \partial P$, let $t_k$ be the first tile of the elimination sequence in $P$. Because no tile in $P \subseteq \{t_k, \dots, t_{71}\}$ contains the reversal of any edge in $L_k$, all lock edges in $L_k$ remain uncancelled on $C$, uniquely forcing $t_k \in P$ and inductively determining $P$.
\end{enumerate}

\section{Algebraic Boundary Framework and Absence of Phasons}
\label{sec:phasons}
In \texttt{AlgebraicBoundary.lean}, we generalize the uniqueness proof to establish that full aperiodic tilings by the $\mathrm{Tile}(1,1)$ Spectre monotile have \textbf{no phasons} (no local finite rearrangements of tiles sharing an identical boundary):
\begin{enumerate}
\item \textbf{Combinatorial Gauss-Bonnet (\texttt{gauss\_bonnet\_patch}):} Every simply-connected finite Spectre patch $P$ has total discrete boundary curvature equal to $12$ step-units ($360^\circ$).
\item \textbf{Transfer-Matrix Curvature Bound (\texttt{lock\_free\_curvature\_bound}):} An exhaustive depth-first search over all punctured 3-gram states (\texttt{verify\_batch1}--\texttt{verify\_batch4} in \texttt{TransferMatrixBatches.lean}) proves that any boundary cycle devoid of the 3 anchor locks has total curvature $\le 2$ step-units ($60^\circ$).
\item \textbf{No Minimal Phasons (\texttt{no\_minimal\_phason}):} Suppose $P_1, P_2$ form a minimal phason ($\partial P_1 = \partial P_2$ with $P_1 \cap P_2 = \emptyset$). By Gauss-Bonnet, their common boundary $w$ has curvature $12$. If $w$ contains a 3-anchor lock, \texttt{lock\_forces\_shared\_tile} forces a shared tile into $P_1 \cap P_2$, a contradiction; if $w$ is lock-free, its curvature is $\le 2$, contradicting $12 \le 2$.
\item \textbf{Global Zero-Phason Theorem (\texttt{no\_phason\_in\_full\_tiling}):} Any hypothetical finite phason flip between two full Spectre tilings $\mathcal{T}_1, \mathcal{T}_2$ reduces, upon stripping shared tiles, to a minimal phason between $\mathcal{T}_1 \setminus \mathcal{T}_2$ and $\mathcal{T}_2 \setminus \mathcal{T}_1$, which is impossible by \texttt{no\_minimal\_phason}.
\end{enumerate}

\section{Kernel Verification Results and Axiom Audit}
\label{sec:results}
All 777 build targets in \texttt{SpectreDeltaBoundary} compile cleanly in the Lean 4 kernel (\texttt{v4.31.0}) with zero warnings, zero errors, zero \texttt{sorryAx}, and zero custom axioms:

\begin{table}[ht]
\centering
\small
\begin{tabular}{ll}
\hline
\textbf{Verified Theorem} & \textbf{Kernel Axiom Footprint} \\
\hline
\texttt{anchor\_lock\_2\_unique}, \texttt{1\_unique}, \texttt{4\_unique} & \texttt{[propext]} \\
\texttt{verify\_delta2\_uniqueness} (71 steps) & \texttt{[propext]} \\
\texttt{verify\_tier1\_gamma\_uniqueness}--\texttt{psi} & \texttt{[propext]} \\
\texttt{delta2\_all\_cascade\_loops\_simple} & \texttt{[propext]} \\
\texttt{delta2\_lock\_elimination\_sequence\_valid} & \texttt{[propext]} \\
\texttt{simple\_closed\_path\_unique\_interior} & \texttt{[propext]} \\
\texttt{no\_minimal\_phason}, \texttt{no\_phason\_in\_full\_tiling} & \texttt{[propext, Classical.choice, Quot.sound]} \\
\texttt{spectre\_hierarchy\_zero\_phasons} & \texttt{[propext, Classical.choice, Quot.sound]} \\
\hline
\end{tabular}
\caption{Lean 4 kernel axiom audit across all four phases of the Aperiodic Holography project.}
\label{tab:perf}
\end{table}

The top-level $\Delta_2$ certificate verification theorem in \texttt{Theorem.lean} is stated as:
\begin{verbatim}
theorem verify_delta2_uniqueness :
  executeCertificate initialMetatileBoundary
    pythonPeelingCertificate = some [] := by
  decide
\end{verbatim}

\section{Discussion}
\label{sec:discussion}
The verified 3-Anchor-Lock boundary reduction cascade and zero-phason theorem demonstrate how localized holographic boundary locks eliminate exponential packing frustration in chiral aperiodic tilings. This architecture provides a general paradigm for computer-assisted proofs of holographic boundary rigidity and zero-phason elasticity in discrete geometry.

\bibliographystyle{plain}
\bibliography{refs}

\end{document}
